% =========================================================
% CAP\'ITULO III: RESULTADO DE LA INVESTIGACI\'ON
% =========================================================

\chapter{Resultado de la investigaci\'on}

\section{Descripci\'on y preparaci\'on del conjunto de datos}

Como fuentes primarias se utilizaron el Sistema Integrado de Estad\'isticas de la Criminalidad y Seguridad Ciudadana del INEI \citep{inei2022sistema} y el Observatorio Nacional de Seguridad Vial del MTC \citep{mtc2024observatorio}, complementados con datos censales y la cartograf\'ia base del INEI para las variables de control. El conjunto final se constituy\'o a partir de los registros de denuncias delictivas y de siniestros de tr\'ansito con informaci\'on geogr\'afica (distrito, comisar\'ia o coordenadas) correspondientes al \'ambito de Lima Metropolitana, cubriendo un periodo de 24 meses consecutivos (julio de 2022 a junio de 2024) con el fin de capturar variaciones estacionales.

En total se recopilaron 212\,486 denuncias por hechos delictivos y 35\,354 siniestros de tr\'ansito. El proceso de preparaci\'on incluy\'o: (i) descarga de los datos publicados en las plataformas oficiales; (ii) normalizaci\'on de los identificadores de distrito, tipo de delito y tipo de siniestro, alineando las taxonom\'ias de ambas fuentes; (iii) geocodificaci\'on de los registros que no contaban con coordenadas, asign\'andoles la localizaci\'on de la comisar\'ia o del punto vial de acuerdo con la cartograf\'ia base; (iv) eliminaci\'on de registros duplicados y de aquellos sin localizaci\'on v\'alida; y (v) agregaci\'on de los eventos individuales en celdas de una malla espacial (\textit{grid}) de 250\,m $\times$ 250\,m y en seis franjas horarias de cuatro horas para el modelado predictivo. Tras la depuraci\'on quedaron 194\,320 denuncias (91.5\,\% del total) y 33\,281 siniestros (94.1\,\% del total), distribuidos en 43\,017 celdas de la malla. La Tabla \ref{tab:datos} resume el conjunto de datos final.

\begin{table}[H]
\centering
\caption[Resumen del conjunto de datos final]{Resumen del conjunto de datos final tras la depuraci\'on (julio 2022--junio 2024, Lima Metropolitana).}
\begin{tabular}{p{4.2cm}p{2.5cm}p{3.2cm}p{2.4cm}p{2.4cm}}
\toprule
Fuente & Dominio & Registros iniciales & Registros finales & Unidad geogr\'afica \\
\midrule
INEI (Sistema Integrado de Estad\'isticas de la Criminalidad) & Denuncias delictivas & 212\,486 & 194\,320 & Comisar\'ia / manzana \\
MTC (Observatorio Nacional de Seguridad Vial) & Siniestros de tr\'ansito & 35\,354 & 33\,281 & Punto / tramo vial \\
\bottomrule
\end{tabular}
\label{tab:datos}
\end{table}

Los delitos m\'as frecuentes fueron el robo y hurto (47.6\,\% del total de denuncias), seguidos del estelionato (12.3\,\%), las lesiones (10.8\,\%), los delitos contra el patrimonio en general (9.4\,\%) y el resto de tipolog\'ias (19.9\,\%). En cuanto a los siniestros de tr\'ansito, la mayor\'ia correspondi\'o a choques (61.2\,\%), seguidos de atropellos (21.4\,\%) y de volcaduras u otro tipo de evento (17.4\,\%). Las estad\'isticas descriptivas b\'asicas (distribuciones por tipo de delito y de siniestro, por distrito, por d\'ia de la semana y por hora) derivan de este an\'alisis y sirven de base para el modelado de las secciones siguientes.

\section{An\'alisis espacial: mapas de calor}

Se elaboraron mapas de calor mediante estimaci\'on de densidad de kernel (KDE) para dos dominios: criminalidad y siniestralidad vial en Lima Metropolitana. En ambos casos se emple\'o un ancho de banda constante de 500\,m, elegido de manera emp\'irica sobre la densidad de puntos, y una malla de celdas de 250\,m, compatible con el nivel de desagregaci\'on de los datos disponibles. Los resultados muestran, de manera consistente con la literatura internacional \citep{chainey2008kde,eck2005mapping}, una concentraci\'on no uniforme del riesgo: las zonas de mayor densidad de denuncias coinciden, en lo sustancial, con los distritos de mayor actividad comercial y densidad poblacional (San Juan de Lurigancho, Comas, San Mart\'in de Porres, Lima Cercado y La Victoria concentran el 38.7\,\% de las denuncias depuradas), mientras que los siniestros de tr\'ansito presentan corredores viales de mayor concentraci\'on asociados a intersecciones de alto flujo y a franjas horarias de mayor congesti\'on vehicular (madrugada y horas de punta). Los principales corredores de siniestralidad identificados fueron la Panamericana Norte y Sur, la avenida Javier Prado, la avenida Abancay, la V\'ia Expresa Luis Bedoya Reyes y la avenida Universitaria.

El an\'alisis temporal mostr\'o patrones diferenciados por dominio: las denuncias por hechos delictivos se concentran entre las 18:00 y las 02:00 horas (48.9\,\% del total), con un pico en la franja de 22:00--02:00 los d\'ias viernes y s\'abado, mientras que los siniestros de tr\'ansito se concentran en las franjas de 22:00--02:00 y 02:00--06:00 (madrugada, 32.6\,\%) y, en segundo orden, en la franja de 06:00--10:00 correspondiente a la hora punta de la ma\~nana (24.1\,\%). Las zonas identificadas como \textit{hotspots} espaciales constituyen la base sobre la que se construye el modelo predictivo del siguiente apartado.

\section{Modelo predictivo espacio-temporal}

Como l\'inea base se consider\'o el modelo de KDE que reproduce la densidad hist\'orica de eventos como estimador de la probabilidad relativa de ocurrencia futura. Sobre esta base se compar\'o un modelo de aprendizaje autom\'atico supervisado --un bosque aleatorio (\textit{random forest})-- que incorpora variables temporales y contextuales (franja horaria, d\'ia de la semana, feriado, distrito, tipo de v\'ia y densidad poblacional de la celda, junto con el valor del mapa de calor KDE como variable adicional). Ambos modelos se entrenaron con un \textit{split} temporal que respeta el orden cronol\'ogico de los datos: entrenamiento con las celdas del periodo julio 2022--diciembre 2023 (18 meses) y evaluaci\'on contra el semestre enero--junio de 2024 (6 meses), no utilizado en el entrenamiento (\textit{backtesting}) \citep{perry2013predictive}.

Para la evaluaci\'on se report\'o el \'Indice de Precisi\'on de Predicci\'on (\textit{Prediction Accuracy Index}, PAI), definido como el cociente entre el porcentaje de eventos futuros capturados en el \'area priorizada y el porcentaje de territorio priorizado, junto con la proporci\'on de territorio que deber\'ia recorrerse para capturar el 50\,\% de los eventos futuros \citep{chainey2008kde}. Los resultados se resumen en la Tabla \ref{tab:comparativo}.

\begin{table}[H]
\centering
\caption[Comparativa de desempe\~no de los modelos evaluados]{Tabla comparativa de desempe\~no de los modelos evaluados (validaci\'on temporal sobre el semestre enero--junio 2024, priorizando el 10\,\% del territorio).}
\begin{tabular}{lcc}
\toprule
Modelo & PAI (10\,\% del territorio priorizado) & Territorio para capturar el 50\,\% de eventos futuros \\
\midrule
KDE (l\'inea base)             & 3.18 & 21.4\,\% \\
Modelo supervisado (bosque aleatorio) & 4.62 & 12.8\,\% \\
\bottomrule
\end{tabular}
\label{tab:comparativo}
\end{table}

Al priorizar el 10\,\% del territorio, la l\'inea base KDE captur\'o el 31.8\,\% de los eventos del periodo de validaci\'on (PAI $=$ 3.18), mientras que el modelo supervisado captur\'o el 46.2\,\% (PAI $=$ 4.62), lo que representa una mejora relativa del 45.3\,\% en eficiencia predictiva espacial. Para capturar la mitad de los eventos futuros, con el modelo KDE habr\'ia que recorrer el 21.4\,\% de la malla, mientras que con el modelo supervisado basta con el 12.8\,\%, lo que se traduce en una reducci\'on del 40.2\,\% del territorio a patrullar manteniendo el mismo nivel de cobertura. Adicionalmente, el modelo supervisado alcanz\'o un AUC-ROC medio de 0.84 (por celda-franja) en la validaci\'on temporal, confirmando su capacidad de discriminaci\'on entre celdas de alto y bajo riesgo. La mejora consistente del modelo supervisado respecto a la l\'inea base KDE justifica su uso como insumo para la planificaci\'on de patrullajes preventivos: las recomendaciones de sectores y franjas horarias priorizados se derivan del ranking de probabilidad generado por este modelo.

\section{Arquitectura de referencia del sistema}

Con el fin de orientar el futuro desarrollo del software, se propuso y document\'o una arquitectura de referencia de tres capas, basada en componentes de c\'odigo abierto: (i) capa de datos, con PostgreSQL y la extensi\'on PostGIS para el almacenamiento y las consultas espaciales \citep{postgis2024}; (ii) capa de procesamiento y modelado, con Python (pandas, geopandas, scikit-learn, scipy) para la limpieza, la construcci\'on de la malla, el c\'alculo de KDE y el entrenamiento de modelos; y (iii) capa de presentaci\'on, con una aplicaci\'on web (Flask o FastAPI en el backend y Leaflet para el mapa interactivo en el frontend) que expone los mapas de calor y las recomendaciones de sectores y horarios. Esta arquitectura fue coherente con el modelo validado en la secci\'on anterior y constituye la gu\'ia para el equipo de desarrollo que contin\'ue el proyecto.

\section{Validaci\'on cualitativa con usuarios potenciales}

La validaci\'on cuantitativa del modelo se efectu\'o mediante simulaci\'on sobre datos hist\'oricos (\textit{backtesting}), tal como se describe en el Cap\'itulo II. La validaci\'on cualitativa con usuarios potenciales (personal de comisar\'ia, gerencia municipal de seguridad ciudadana o \'area de suscripci\'on de riesgos de una aseguradora) se deja como l\'inea de trabajo futura, dada la naturaleza experimental del presente proyecto. Esta limitaci\'on se reconoce expl\'icitamente como una de las razones por las que los resultados de este cap\'itulo constituyen evidencia de viabilidad t\'ecnica y no una evaluaci\'on operativa completa.

\section{Discusi\'on de resultados}

Los resultados del Cap\'itulo III muestran que un modelo de anal\'itica predictiva espacio-temporal construido sobre datos abiertos permite estimar, con un desempe\~no superior a la l\'inea base KDE, las zonas y franjas horarias de mayor riesgo de criminalidad y siniestralidad vial en Lima Metropolitana. Esto es consistente con la evidencia internacional revisada en el Cap\'itulo II \citep{chainey2008kde,mohler2011selfexciting,perry2013predictive}, y con los antecedentes nacionales que identifican la fragmentaci\'on de fuentes y la carencia de un enfoque predictivo como brechas operativas \citep{sanchezguevara2022tabcop,mora2015sistematizacion}. Los patrones horarios de siniestralidad observados en este trabajo coinciden, en lo sustancial, con lo reportado por el MTC sobre la mayor concentraci\'on de siniestros en horas de baja visibilidad y en franjas de mayor congesti\'on vehicular \citep{mtc2024observatorio}. Para la planificaci\'on del patrullaje, las implicancias pr\'acticas son que las recomendaciones de sectores y horarios generadas por el modelo pueden servir de insumo objetivo para complementar el criterio experto del personal de patrullaje, siempre dentro del nivel agregado espacio-temporal y respetando las salvaguardas \'eticas discutidas en el Cap\'itulo II.